% =============================================================================
\section{Introducción}
% =============================================================================

El estudio del rebote inelástico de partículas y clústeres nanométricos sobre
superficies cristalinas tiene relevancia directa en procesos de deposición de
capas delgadas~\cite{Yamaguchi1992}, erosión superficial por impacto~\cite{Zimmermann2021},
síntesis de nanopartículas~\cite{Kluge2020} y tribología a escala
atómica~\cite{Szlufarska2008}. A escala nanométrica, los efectos cuánticos son
despreciables pero los efectos térmicos y las fuerzas de adhesión superficial
son comparables a la energía cinética del proyectil, dando lugar a regímenes de
comportamiento cualitativamente distintos del caso macroscópico.

La dinámica molecular (MD) es la herramienta natural para simular estos procesos:
integra las ecuaciones de Newton para sistemas de $10^2$--$10^5$ átomos en
tiempos de $10^{1}$--$10^3$~ps con potenciales empíricos bien validados para
gases nobles~\cite{Lennard-Jones1931}. Sin embargo, el salto de ``correr una
simulación'' a ``producir resultados publicables y reproducibles'' involucra
desafíos de ingeniería de software que rara vez se discuten en la literatura
científica: gestión de parámetros, validación de resultados, barridos
paramétricos paralelos, caché de análisis y reproducibilidad del entorno.

\subsection{Motivación pedagógica}

Este trabajo nace en el contexto de un Magíster en Simulación Computacional
(PUCV, Chile), donde uno de los objetivos es demostrar que los patrones de
\textbf{programación avanzada} ---metaclases, descriptores, decoradores,
generadores, jerarquías de excepciones y análisis tabular con
\texttt{pandas}--- no son un fin en sí mismos, sino herramientas que resuelven
problemas concretos y mensurables en investigación científica computacional.

La pregunta pedagógica central es: \textit{¿qué aporta cada patrón de
programación avanzada al flujo de trabajo de un investigador en MD?} Para
responderla, se documentan explícitamente los problemas que cada patrón resuelve
y la ganancia observada en reproducibilidad, mantenibilidad y eficiencia.

\subsection{Contribuciones del trabajo}

Este artículo presenta las siguientes contribuciones:

\begin{enumerate}
    \item Un \textit{pipeline} completo de MD en \python{}~3.12 +
          \lammps{}~7~Feb~2024 para el estudio de rebote inelástico de
          clústeres LJ, arquitecturado con patrones de programación avanzada
          (Sección~\ref{sec:arquitectura}).

    \item Resultados físicos de la dependencia $\cor(N, v_0, \Tstar, \theta,
          \epsps)$ y la partición energética $\alphaint(N)$ para clústeres de
          tamaños mágicos icosaédricos $N \in \{1, 2, 7, 13, 55\}$
          (Sección~\ref{sec:resultados}).

    \item La introducción del número de Weber $\We = \Einc/(N\epsps)$ como
          variable de escala que colapsa las curvas de transición de adhesión
          para distintos $N$ y $\epsps$ (Sección~\ref{sec:resultados}).

    \item Una discusión explícita del valor pedagógico de cada decisión de
          diseño del software (Sección~\ref{sec:discusion}).
\end{enumerate}

El código fuente completo está disponible
en \url{https://github.com/fwasaff/inelastic-rebound-md} bajo licencia MIT,
incluyendo scripts de \lammps{}, todos los módulos Python, conjuntos de datos
procesados y este documento.

\subsection{Trabajos relacionados}

El rebote inelástico átomo-superficie ha sido ampliamente estudiado con MD para
sistemas LJ~\cite{Tully1980, Head-Gordon1991, Binnie2004} y con potenciales más
realistas (EAM, ReaxFF) para metales~\cite{Foiles1986}. Para clústeres, los
trabajos de Zimmermann~et~al.~\cite{Zimmermann2021} y Kluge~et~al.~\cite{Kluge2020}
han estudiado deposición sobre superficies metálicas, si bien con énfasis en
morfología de depósito más que en energética del rebote. El número de Weber como
parámetro de escala fue introducido por Dahneke~\cite{Dahneke1975} para
aerosoles y posteriormente extendido a nanopartículas~\cite{Hou2000,Inaba2004}.

En cuanto a la reproducibilidad computacional en MD, iniciativas como
\textsc{Moltemplate}~\cite{Jewett2021} y \textsc{pyiron}~\cite{Janssen2019}
abordan la gestión de flujos de trabajo, pero con énfasis en sistemas más
complejos y menor énfasis en la pedagogía del diseño de software.
